\documentclass[12pt]{article}
\usepackage{amsmath}
\usepackage{amssymb}
\usepackage{geometry}
\usepackage{xcolor}
\geometry{margin=1in}

\title{GED Mathematics Worksheet - Rational Expressions and Functions -- Full Solutions}
%\subtitle{Rational Expressions and Functions -- Full Solutions}
\author{}
\date{}

\begin{document}

\maketitle

\section*{Part 1: Multiplying Rational Expressions (Questions 1-10)}

\subsection*{Question 1}
\textbf{Problem:} Multiply: $\frac{3x}{4} \cdot \frac{8}{9x}$

\textbf{Solution:}
\begin{align*}
\frac{3x}{4} \cdot \frac{8}{9x} &= \frac{3x \cdot 8}{4 \cdot 9x} \\
&= \frac{24x}{36x} \\
&= \frac{24}{36} \quad \text{(cancel } x \text{)} \\
&= \frac{2}{3} \quad \text{(simplify by dividing by 12)}
\end{align*}

\textbf{Answer:} $\boxed{\frac{2}{3}}$

\subsection*{Question 2}
\textbf{Problem:} Multiply: $\frac{5a^2}{6b} \cdot \frac{12b}{10a}$

\textbf{Solution:}
\begin{align*}
\frac{5a^2}{6b} \cdot \frac{12b}{10a} &= \frac{5a^2 \cdot 12b}{6b \cdot 10a} \\
&= \frac{60a^2b}{60ab} \\
&= \frac{a^2}{a} \quad \text{(cancel } 60b \text{)} \\
&= a
\end{align*}

\textbf{Answer:} $\boxed{a}$

\subsection*{Question 3}
\textbf{Problem:} Multiply: $\frac{x+2}{x-1} \cdot \frac{x-1}{x+3}$

\textbf{Solution:}
\begin{align*}
\frac{x+2}{x-1} \cdot \frac{x-1}{x+3} &= \frac{(x+2)(x-1)}{(x-1)(x+3)} \\
&= \frac{x+2}{x+3} \quad \text{(cancel } (x-1) \text{)} \\
\end{align*}

\textbf{Restrictions:} $x \neq 1, -3$

\textbf{Answer:} $\boxed{\frac{x+2}{x+3}}$

\subsection*{Question 4}
\textbf{Problem:} Multiply: $\frac{2x^2}{5y} \cdot \frac{15y^2}{8x}$

\textbf{Solution:}
\begin{align*}
\frac{2x^2}{5y} \cdot \frac{15y^2}{8x} &= \frac{2x^2 \cdot 15y^2}{5y \cdot 8x} \\
&= \frac{30x^2y^2}{40xy} \\
&= \frac{30xy}{40} \quad \text{(cancel } x \text{ and } y \text{)} \\
&= \frac{3xy}{4} \quad \text{(divide by 10)}
\end{align*}

\textbf{Answer:} $\boxed{\frac{3xy}{4}}$

\subsection*{Question 5}
\textbf{Problem:} Multiply: $\frac{(x+4)(x-3)}{(x+1)} \cdot \frac{(x+5)}{(x-3)}$

\textbf{Solution:}
\begin{align*}
\frac{(x+4)(x-3)}{(x+1)} \cdot \frac{(x+5)}{(x-3)} &= \frac{(x+4)(x-3)(x+5)}{(x+1)(x-3)} \\
&= \frac{(x+4)(x+5)}{(x+1)} \quad \text{(cancel } (x-3) \text{)} \\
&= \frac{x^2 + 5x + 4x + 20}{x+1} \quad \text{(FOIL)} \\
&= \frac{x^2 + 9x + 20}{x+1}
\end{align*}

\textbf{Restrictions:} $x \neq -1, 3$

\textbf{Answer:} $\boxed{\frac{(x+4)(x+5)}{x+1} \text{ or } \frac{x^2+9x+20}{x+1}}$

\subsection*{Question 6}
\textbf{Problem:} Multiply: $\frac{x^2-4}{x^2+5x+6} \cdot \frac{x+3}{x-2}$

\textbf{Solution:}

First, factor all numerators and denominators:
\begin{align*}
x^2 - 4 &= (x-2)(x+2) \\
x^2 + 5x + 6 &= (x+2)(x+3)
\end{align*}

Now multiply:
\begin{align*}
\frac{(x-2)(x+2)}{(x+2)(x+3)} \cdot \frac{x+3}{x-2} &= \frac{(x-2)(x+2)(x+3)}{(x+2)(x+3)(x-2)} \\
&= 1 \quad \text{(cancel all factors)}
\end{align*}

\textbf{Restrictions:} $x \neq -3, -2, 2$

\textbf{Answer:} $\boxed{1}$

\subsection*{Question 7}
\textbf{Problem:} Multiply: $\frac{6m^2n}{7} \cdot \frac{14}{18mn^2}$

\textbf{Solution:}
\begin{align*}
\frac{6m^2n}{7} \cdot \frac{14}{18mn^2} &= \frac{6m^2n \cdot 14}{7 \cdot 18mn^2} \\
&= \frac{84m^2n}{126mn^2} \\
&= \frac{84m}{126n} \quad \text{(cancel } n \text{)} \\
&= \frac{4m}{6n} \quad \text{(divide by 21)} \\
&= \frac{2m}{3n} \quad \text{(divide by 2)}
\end{align*}

\textbf{Answer:} $\boxed{\frac{2m}{3n}}$

\subsection*{Question 8}
\textbf{Problem:} Multiply: $\frac{x^2-9}{x^2-1} \cdot \frac{x+1}{x+3}$

\textbf{Solution:}

Factor:
\begin{align*}
x^2 - 9 &= (x-3)(x+3) \\
x^2 - 1 &= (x-1)(x+1)
\end{align*}

Multiply:
\begin{align*}
\frac{(x-3)(x+3)}{(x-1)(x+1)} \cdot \frac{x+1}{x+3} &= \frac{(x-3)(x+3)(x+1)}{(x-1)(x+1)(x+3)} \\
&= \frac{x-3}{x-1} \quad \text{(cancel } (x+3) \text{ and } (x+1) \text{)}
\end{align*}

\textbf{Restrictions:} $x \neq 1, -1, -3$

\textbf{Answer:} $\boxed{\frac{x-3}{x-1}}$

\subsection*{Question 9}
\textbf{Problem:} Multiply: $\frac{a^2-a-6}{a^2+4a+3} \cdot \frac{a+1}{a-3}$

\textbf{Solution:}

Factor:
\begin{align*}
a^2 - a - 6 &= (a-3)(a+2) \quad \text{(factors of } -6 \text{ that differ by } -1 \text{: } -3, 2 \text{)} \\
a^2 + 4a + 3 &= (a+1)(a+3) \quad \text{(factors of } 3 \text{ that sum to } 4 \text{: } 1, 3 \text{)}
\end{align*}

Multiply:
\begin{align*}
\frac{(a-3)(a+2)}{(a+1)(a+3)} \cdot \frac{a+1}{a-3} &= \frac{(a-3)(a+2)(a+1)}{(a+1)(a+3)(a-3)} \\
&= \frac{a+2}{a+3} \quad \text{(cancel } (a-3) \text{ and } (a+1) \text{)}
\end{align*}

\textbf{Restrictions:} $a \neq -1, -3, 3$

\textbf{Answer:} $\boxed{\frac{a+2}{a+3}}$

\subsection*{Question 10}
\textbf{Problem:} Multiply: $\frac{2x^2-8x}{x^2-16} \cdot \frac{x+4}{2x}$

\textbf{Solution:}

Factor:
\begin{align*}
2x^2 - 8x &= 2x(x-4) \\
x^2 - 16 &= (x-4)(x+4)
\end{align*}

Multiply:
\begin{align*}
\frac{2x(x-4)}{(x-4)(x+4)} \cdot \frac{x+4}{2x} &= \frac{2x(x-4)(x+4)}{(x-4)(x+4)(2x)} \\
&= 1 \quad \text{(cancel all factors)}
\end{align*}

\textbf{Restrictions:} $x \neq 0, 4, -4$

\textbf{Answer:} $\boxed{1}$

\newpage

\section*{Part 2: Dividing Rational Expressions (Questions 11-20)}

\subsection*{Question 11}
\textbf{Problem:} Divide: $\frac{5}{8} \div \frac{3}{4}$

\textbf{Solution:}
\begin{align*}
\frac{5}{8} \div \frac{3}{4} &= \frac{5}{8} \cdot \frac{4}{3} \quad \text{(invert and multiply)} \\
&= \frac{5 \cdot 4}{8 \cdot 3} \\
&= \frac{20}{24} \\
&= \frac{5}{6} \quad \text{(divide by 4)}
\end{align*}

\textbf{Answer:} $\boxed{\frac{5}{6}}$

\subsection*{Question 12}
\textbf{Problem:} Divide: $\frac{x^2}{y} \div \frac{x}{2y}$

\textbf{Solution:}
\begin{align*}
\frac{x^2}{y} \div \frac{x}{2y} &= \frac{x^2}{y} \cdot \frac{2y}{x} \quad \text{(invert and multiply)} \\
&= \frac{x^2 \cdot 2y}{y \cdot x} \\
&= \frac{2x^2y}{xy} \\
&= 2x \quad \text{(cancel } xy \text{)}
\end{align*}

\textbf{Answer:} $\boxed{2x}$

\subsection*{Question 13}
\textbf{Problem:} Divide: $\frac{3a}{5b} \div \frac{9a^2}{10b^2}$

\textbf{Solution:}
\begin{align*}
\frac{3a}{5b} \div \frac{9a^2}{10b^2} &= \frac{3a}{5b} \cdot \frac{10b^2}{9a^2} \quad \text{(invert and multiply)} \\
&= \frac{3a \cdot 10b^2}{5b \cdot 9a^2} \\
&= \frac{30ab^2}{45a^2b} \\
&= \frac{30b}{45a} \quad \text{(cancel } ab \text{)} \\
&= \frac{2b}{3a} \quad \text{(divide by 15)}
\end{align*}

\textbf{Answer:} $\boxed{\frac{2b}{3a}}$

\subsection*{Question 14}
\textbf{Problem:} Divide: $\frac{m^2n}{6} \div \frac{mn^2}{12}$

\textbf{Solution:}
\begin{align*}
\frac{m^2n}{6} \div \frac{mn^2}{12} &= \frac{m^2n}{6} \cdot \frac{12}{mn^2} \quad \text{(invert and multiply)} \\
&= \frac{m^2n \cdot 12}{6 \cdot mn^2} \\
&= \frac{12m^2n}{6mn^2} \\
&= \frac{12m}{6n} \quad \text{(cancel } mn \text{)} \\
&= 2m \quad \text{(divide by 6)}
\end{align*}

\textbf{Answer:} $\boxed{2m}$

\subsection*{Question 15}
\textbf{Problem:} Divide: $\frac{x+2}{x-3} \div \frac{x+2}{x+5}$

\textbf{Solution:}
\begin{align*}
\frac{x+2}{x-3} \div \frac{x+2}{x+5} &= \frac{x+2}{x-3} \cdot \frac{x+5}{x+2} \quad \text{(invert and multiply)} \\
&= \frac{(x+2)(x+5)}{(x-3)(x+2)} \\
&= \frac{x+5}{x-3} \quad \text{(cancel } (x+2) \text{)}
\end{align*}

\textbf{Restrictions:} $x \neq -2, 3, -5$

\textbf{Answer:} $\boxed{\frac{x+5}{x-3}}$

\subsection*{Question 16}
\textbf{Problem:} Divide: $\frac{x^2-4}{x^2+6x+9} \div \frac{x-2}{x+3}$

\textbf{Solution:}

Factor:
\begin{align*}
x^2 - 4 &= (x-2)(x+2) \\
x^2 + 6x + 9 &= (x+3)^2
\end{align*}

Divide:
\begin{align*}
\frac{(x-2)(x+2)}{(x+3)^2} \div \frac{x-2}{x+3} &= \frac{(x-2)(x+2)}{(x+3)^2} \cdot \frac{x+3}{x-2} \quad \text{(invert and multiply)} \\
&= \frac{(x-2)(x+2)(x+3)}{(x+3)^2(x-2)} \\
&= \frac{x+2}{x+3} \quad \text{(cancel } (x-2) \text{ and one } (x+3) \text{)}
\end{align*}

\textbf{Restrictions:} $x \neq -3, 2$

\textbf{Answer:} $\boxed{\frac{x+2}{x+3}}$

\subsection*{Question 17}
\textbf{Problem:} Divide: $\frac{x^2-8x}{x^2+8x+16} \div \frac{x-8}{x^2-2x-24}$

\textbf{Solution:}

Factor:
\begin{align*}
x^2 - 8x &= x(x-8) \\
x^2 + 8x + 16 &= (x+4)^2 \\
x^2 - 2x - 24 &= (x-6)(x+4) \quad \text{(factors of } -24 \text{ that differ by } -2 \text{: } -6, 4 \text{)}
\end{align*}

Divide:
\begin{align*}
\frac{x(x-8)}{(x+4)^2} \div \frac{x-8}{(x-6)(x+4)} &= \frac{x(x-8)}{(x+4)^2} \cdot \frac{(x-6)(x+4)}{x-8} \quad \text{(invert and multiply)} \\
&= \frac{x(x-8)(x-6)(x+4)}{(x+4)^2(x-8)} \\
&= \frac{x(x-6)}{x+4} \quad \text{(cancel } (x-8) \text{ and one } (x+4) \text{)}
\end{align*}

\textbf{Restrictions:} $x \neq 0, -4, 8, 6$

\textbf{Answer:} $\boxed{\frac{x(x-6)}{x+4}}$

\subsection*{Question 18}
\textbf{Problem:} Divide: $\frac{2x^2+5x+3}{x^2-1} \div \frac{2x+3}{x-1}$

\textbf{Solution:}

Factor:
\begin{align*}
2x^2 + 5x + 3 &= (2x+3)(x+1) \quad \text{(factors of } 6 \text{ that sum to } 5 \text{: } 2, 3 \text{)} \\
x^2 - 1 &= (x-1)(x+1)
\end{align*}

Divide:
\begin{align*}
\frac{(2x+3)(x+1)}{(x-1)(x+1)} \div \frac{2x+3}{x-1} &= \frac{(2x+3)(x+1)}{(x-1)(x+1)} \cdot \frac{x-1}{2x+3} \quad \text{(invert and multiply)} \\
&= \frac{(2x+3)(x+1)(x-1)}{(x-1)(x+1)(2x+3)} \\
&= 1 \quad \text{(cancel all factors)}
\end{align*}

\textbf{Restrictions:} $x \neq 1, -1, -\frac{3}{2}$

\textbf{Answer:} $\boxed{1}$

\subsection*{Question 19}
\textbf{Problem:} Divide: $\frac{a^2-9}{a^2-4} \div \frac{a+3}{a-2}$

\textbf{Solution:}

Factor:
\begin{align*}
a^2 - 9 &= (a-3)(a+3) \\
a^2 - 4 &= (a-2)(a+2)
\end{align*}

Divide:
\begin{align*}
\frac{(a-3)(a+3)}{(a-2)(a+2)} \div \frac{a+3}{a-2} &= \frac{(a-3)(a+3)}{(a-2)(a+2)} \cdot \frac{a-2}{a+3} \quad \text{(invert and multiply)} \\
&= \frac{(a-3)(a+3)(a-2)}{(a-2)(a+2)(a+3)} \\
&= \frac{a-3}{a+2} \quad \text{(cancel } (a+3) \text{ and } (a-2) \text{)}
\end{align*}

\textbf{Restrictions:} $a \neq 2, -2, -3$

\textbf{Answer:} $\boxed{\frac{a-3}{a+2}}$

\subsection*{Question 20}
\textbf{Problem:} Divide: $\frac{x^2-5x+6}{x^2+x-12} \div \frac{x-2}{x+4}$

\textbf{Solution:}

Factor:
\begin{align*}
x^2 - 5x + 6 &= (x-2)(x-3) \quad \text{(factors of } 6 \text{ that sum to } -5 \text{: } -2, -3 \text{)} \\
x^2 + x - 12 &= (x+4)(x-3) \quad \text{(factors of } -12 \text{ that sum to } 1 \text{: } 4, -3 \text{)}
\end{align*}

Divide:
\begin{align*}
\frac{(x-2)(x-3)}{(x+4)(x-3)} \div \frac{x-2}{x+4} &= \frac{(x-2)(x-3)}{(x+4)(x-3)} \cdot \frac{x+4}{x-2} \quad \text{(invert and multiply)} \\
&= \frac{(x-2)(x-3)(x+4)}{(x+4)(x-3)(x-2)} \\
&= 1 \quad \text{(cancel all factors)}
\end{align*}

\textbf{Restrictions:} $x \neq -4, 3, 2$

\textbf{Answer:} $\boxed{1}$

\newpage

\section*{Part 3: Understanding Functions (Questions 21-30)}

\subsection*{Question 21}
\textbf{Problem:} Determine if the relation $\{(1, 2), (2, 4), (3, 6), (4, 8)\}$ is a function. Explain why or why not.

\textbf{Solution:}

A relation is a function if each element in the domain is paired with exactly one element in the range.

Domain: $\{1, 2, 3, 4\}$
Range: $\{2, 4, 6, 8\}$

Checking each domain element:
\begin{itemize}
\item $1 \to 2$ (one output)
\item $2 \to 4$ (one output)
\item $3 \to 6$ (one output)
\item $4 \to 8$ (one output)
\end{itemize}

Each domain element is paired with exactly one range element.

\textbf{Answer:} $\boxed{\text{Yes, it is a function.}}$

\subsection*{Question 22}
\textbf{Problem:} Determine if the relation $\{(1, 3), (2, 5), (1, 7), (3, 9)\}$ is a function. Explain why or why not.

\textbf{Solution:}

Domain: $\{1, 2, 3\}$
Range: $\{3, 5, 7, 9\}$

Checking each domain element:
\begin{itemize}
\item $1 \to 3$ and $1 \to 7$ (two outputs!)
\item $2 \to 5$ (one output)
\item $3 \to 9$ (one output)
\end{itemize}

The domain element $1$ is paired with two different range elements ($3$ and $7$).

\textbf{Answer:} $\boxed{\text{No, it is not a function.}}$

\subsection*{Question 23}
\textbf{Problem:} For $f(x) = 3x - 2$, find $f(5)$.

\textbf{Solution:}

Substitute $x = 5$ into the function:
\begin{align*}
f(5) &= 3(5) - 2 \\
&= 15 - 2 \\
&= 13
\end{align*}

\textbf{Answer:} $\boxed{f(5) = 13}$

\subsection*{Question 24}
\textbf{Problem:} For $g(x) = x^2 + 4$, find $g(-3)$.

\textbf{Solution:}

Substitute $x = -3$ into the function:
\begin{align*}
g(-3) &= (-3)^2 + 4 \\
&= 9 + 4 \\
&= 13
\end{align*}

\textbf{Answer:} $\boxed{g(-3) = 13}$

\subsection*{Question 25}
\textbf{Problem:} For $h(x) = 2x^2 - 5x + 1$, find $h(2)$.

\textbf{Solution:}

Substitute $x = 2$ into the function:
\begin{align*}
h(2) &= 2(2)^2 - 5(2) + 1 \\
&= 2(4) - 10 + 1 \\
&= 8 - 10 + 1 \\
&= -1
\end{align*}

\textbf{Answer:} $\boxed{h(2) = -1}$

\subsection*{Question 26}
\textbf{Problem:} For $f(x) = \frac{x+3}{x-1}$, find $f(4)$.

\textbf{Solution:}

Substitute $x = 4$ into the function:
\begin{align*}
f(4) &= \frac{4+3}{4-1} \\
&= \frac{7}{3}
\end{align*}

\textbf{Answer:} $\boxed{f(4) = \frac{7}{3}}$

\subsection*{Question 27}
\textbf{Problem:} For $g(x) = 5x - 7$, find $g(-2)$.

\textbf{Solution:}

Substitute $x = -2$ into the function:
\begin{align*}
g(-2) &= 5(-2) - 7 \\
&= -10 - 7 \\
&= -17
\end{align*}

\textbf{Answer:} $\boxed{g(-2) = -17}$

\subsection*{Question 28}
\textbf{Problem:} For $h(x) = x^3 - 2x + 4$, find $h(3)$.

\textbf{Solution:}

Substitute $x = 3$ into the function:
\begin{align*}
h(3) &= (3)^3 - 2(3) + 4 \\
&= 27 - 6 + 4 \\
&= 25
\end{align*}

\textbf{Answer:} $\boxed{h(3) = 25}$

\subsection*{Question 29}
\textbf{Problem:} For $f(x) = \sqrt{x + 5}$, find $f(4)$.

\textbf{Solution:}

Substitute $x = 4$ into the function:
\begin{align*}
f(4) &= \sqrt{4 + 5} \\
&= \sqrt{9} \\
&= 3
\end{align*}

\textbf{Answer:} $\boxed{f(4) = 3}$

\subsection*{Question 30}
\textbf{Problem:} For $g(x) = |2x - 3|$, find $g(1)$.

\textbf{Solution:}

Substitute $x = 1$ into the function:
\begin{align*}
g(1) &= |2(1) - 3| \\
&= |2 - 3| \\
&= |-1| \\
&= 1
\end{align*}

\textbf{Answer:} $\boxed{g(1) = 1}$

\newpage

\section*{Part 4: Evaluating Functions with Expressions (Questions 31-35)}

\subsection*{Question 31}
\textbf{Problem:} For $g(x) = 3x + 5$, find $g(4n)$.

\textbf{Solution:}

Substitute $x = 4n$ into the function:
\begin{align*}
g(4n) &= 3(4n) + 5 \\
&= 12n + 5
\end{align*}

\textbf{Answer:} $\boxed{g(4n) = 12n + 5}$

\subsection*{Question 32}
\textbf{Problem:} For $h(x) = 2x^2 - 3x + 7$, find $h(4y^2)$.

\textbf{Solution:}

Substitute $x = 4y^2$ into the function:
\begin{align*}
h(4y^2) &= 2(4y^2)^2 - 3(4y^2) + 7 \\
&= 2(16y^4) - 12y^2 + 7 \\
&= 32y^4 - 12y^2 + 7
\end{align*}

\textbf{Answer:} $\boxed{h(4y^2) = 32y^4 - 12y^2 + 7}$

\subsection*{Question 33}
\textbf{Problem:} For $f(x) = x^2 - 6x$, find $f(a+2)$.

\textbf{Solution:}

Substitute $x = a+2$ into the function:
\begin{align*}
f(a+2) &= (a+2)^2 - 6(a+2) \\
&= a^2 + 4a + 4 - 6a - 12 \\
&= a^2 - 2a - 8
\end{align*}

\textbf{Answer:} $\boxed{f(a+2) = a^2 - 2a - 8}$

\subsection*{Question 34}
\textbf{Problem:} For $g(x) = 5x - 1$, find $g(2m + 3)$.

\textbf{Solution:}

Substitute $x = 2m + 3$ into the function:
\begin{align*}
g(2m+3) &= 5(2m+3) - 1 \\
&= 10m + 15 - 1 \\
&= 10m + 14
\end{align*}

\textbf{Answer:} $\boxed{g(2m+3) = 10m + 14}$

\subsection*{Question 35}
\textbf{Problem:} For $h(x) = x^2 + 2x - 8$, find $h(x-1)$.

\textbf{Solution:}

Substitute $x = x-1$ into the function:
\begin{align*}
h(x-1) &= (x-1)^2 + 2(x-1) - 8 \\
&= x^2 - 2x + 1 + 2x - 2 - 8 \\
&= x^2 - 9
\end{align*}

Wait, let me recalculate more carefully:
\begin{align*}
h(x-1) &= (x-1)^2 + 2(x-1) - 8 \\
&= (x^2 - 2x + 1) + (2x - 2) - 8 \\
&= x^2 - 2x + 1 + 2x - 2 - 8 \\
&= x^2 + (- 2x + 2x) + (1 - 2 - 8) \\
&= x^2 - 9
\end{align*}

Actually, let me verify once more by expanding step-by-step:
\begin{align*}
h(x-1) &= (x-1)^2 + 2(x-1) - 8 \\
&= x^2 - 2x + 1 + 2x - 2 - 8 \\
&= x^2 + (-2x + 2x) + (1 - 2 - 8) \\
&= x^2 + 0 + (-9) \\
&= x^2 - 9
\end{align*}

Hmm, but this doesn't match the answer key. Let me reconsider. The original function is $h(x) = x^2 + 2x - 8$. When we substitute $x-1$:

\begin{align*}
h(x-1) &= (x-1)^2 + 2(x-1) - 8 \\
&= x^2 - 2x + 1 + 2x - 2 - 8 \\
&= x^2 - 9
\end{align*}

Wait, I think the answer key might have a different interpretation. Let me check if they meant to find $h(t)$ where $t = x - 1$, then express in terms of $x$. Actually, the standard interpretation is what I did. But the answer key says $x^2 - 4x - 5$. Let me recalculate:

Actually, I realize my error. Let me be more careful:
\begin{align*}
h(x-1) &= (x-1)^2 + 2(x-1) - 8
\end{align*}

Let $u = x - 1$. Then:
\begin{align*}
h(u) &= u^2 + 2u - 8
\end{align*}

So:
\begin{align*}
h(x-1) &= (x-1)^2 + 2(x-1) - 8 \\
&= x^2 - 2x + 1 + 2x - 2 - 8 \\
&= x^2 - 9
\end{align*}

But the answer key says $x^2 - 4x - 5$. Let me check if there's a different function. Actually, I think there might be an error in my answer key. Let me verify by a different method. If $h(x) = x^2 + 2x - 8$, then:

$h(x-1) = (x-1)^2 + 2(x-1) - 8 = x^2 - 2x + 1 + 2x - 2 - 8 = x^2 - 9$

This is correct. However, the provided answer key says $x^2 - 4x - 5$. Let me check if perhaps the original function was different. Looking back at the worksheet, it says $h(x) = x^2 + 2x - 8$. 

Actually, I wonder if the answer key had a typo. Let me just provide the correct solution:

\textbf{Answer:} $\boxed{h(x-1) = x^2 - 9}$

Actually, wait. Let me reconsider the answer key. Perhaps I made an error. Let me expand very carefully:

$h(x-1) = (x-1)^2 + 2(x-1) - 8$

$(x-1)^2 = x^2 - 2x + 1$
$2(x-1) = 2x - 2$

So: $h(x-1) = x^2 - 2x + 1 + 2x - 2 - 8 = x^2 + (-2x + 2x) + (1 - 2 - 8) = x^2 - 9$

I'm confident this is correct. The answer key might have had a different problem in mind. Let me provide the correct answer:

\textbf{Answer:} $\boxed{h(x-1) = x^2 - 9}$

\newpage

\section*{Part 5: Functions and Graphs (Questions 36-40)}

\subsection*{Question 36}
\textbf{Problem:} Use the vertical line test to determine if the graph of $y = x^2$ represents a function.

\textbf{Solution:}

The vertical line test states: A relation is a function if every vertical line intersects the graph at most once.

For $y = x^2$ (a parabola opening upward):
\begin{itemize}
\item The parabola is symmetric about the $y$-axis
\item For any value of $x$, there is exactly one corresponding $y$-value
\item Any vertical line drawn on the graph will intersect the parabola at exactly one point
\end{itemize}

Therefore, the graph passes the vertical line test.

\textbf{Answer:} $\boxed{\text{Yes, } y = x^2 \text{ represents a function.}}$

\subsection*{Question 37}
\textbf{Problem:} Use the vertical line test to determine if the graph of $x = y^2$ represents a function.

\textbf{Solution:}

For $x = y^2$ (a parabola opening to the right):
\begin{itemize}
\item This can be rewritten as $y = \pm\sqrt{x}$
\item For any positive value of $x$, there are two corresponding $y$-values (one positive and one negative)
\item For example, when $x = 4$: $y = 2$ or $y = -2$
\item A vertical line at $x = 4$ would intersect the parabola at two points: $(4, 2)$ and $(4, -2)$
\end{itemize}

Therefore, the graph fails the vertical line test.

\textbf{Answer:} $\boxed{\text{No, } x = y^2 \text{ does not represent a function.}}$

\subsection*{Question 38}
\textbf{Problem:} For the function $f(x) = 2x - 3$, find $f(2)$ using the equation. Then verify your answer by substituting into the function.

\textbf{Solution:}

Using the equation:
\begin{align*}
f(2) &= 2(2) - 3 \\
&= 4 - 3 \\
&= 1
\end{align*}

Verification by substitution:
\begin{align*}
f(x) &= 2x - 3 \\
f(2) &= 2(2) - 3 = 4 - 3 = 1 \quad \checkmark
\end{align*}

Both methods give the same result.

\textbf{Answer:} $\boxed{f(2) = 1}$

\subsection*{Question 39}
\textbf{Problem:} Describe the behavior of a function during the interval $-2 \leq x < 1$ if the function increases throughout this interval.

\textbf{Solution:}

If a function increases throughout the interval $-2 \leq x < 1$, this means:

\begin{itemize}
\item As $x$ increases from $-2$ to values approaching $1$, the $y$-values (function outputs) also increase
\item The function has a positive slope or positive rate of change in this interval
\item If we pick any two points $x_1$ and $x_2$ where $-2 \leq x_1 < x_2 < 1$, then $f(x_1) < f(x_2)$
\item Graphically, the function rises from left to right throughout this interval
\item The function reaches its minimum value at $x = -2$ and approaches its maximum value as $x$ approaches $1$ (but doesn't reach it since the interval is open at $x = 1$)
\end{itemize}

\textbf{Answer:} $\boxed{\text{The function increases as } x \text{ increases; the graph rises from left to right.}}$

\subsection*{Question 40}
\textbf{Problem:} For a piecewise function, explain why different intervals might have different behaviors (increasing, decreasing, or constant).

\textbf{Solution:}

A piecewise function is defined by different rules (equations) on different intervals of the domain. Because each interval uses a different rule, the behavior can vary:

\begin{itemize}
\item \textbf{Different Rules:} Each piece of the function follows a different equation. For example:
  \begin{align*}
  f(x) = \begin{cases}
  x + 2 & \text{if } x < 0 \\
  x^2 & \text{if } 0 \leq x < 3 \\
  5 & \text{if } x \geq 3
  \end{cases}
  \end{align*}

\item \textbf{Linear vs. Nonlinear:} One piece might be linear (constant slope), while another is quadratic (changing slope)

\item \textbf{Different Slopes:} Linear pieces can have different slopes. For example, $f(x) = x + 2$ has slope $1$, while $f(x) = -2x$ has slope $-2$

\item \textbf{Behavior Examples:}
  \begin{itemize}
  \item For $x < 0$: $f(x) = x + 2$ increases (positive slope)
  \item For $0 \leq x < 3$: $f(x) = x^2$ increases (parabola opening upward)
  \item For $x \geq 3$: $f(x) = 5$ is constant (horizontal line)
  \end{itemize}

\item \textbf{Transitions:} At the boundaries between intervals, the function can have jumps, corners, or smooth transitions depending on how the pieces connect
\end{itemize}

\textbf{Answer:} $\boxed{\text{Different intervals use different rules, so each piece can have its own behavior.}}$

\end{document}